We propose several directions for future research. In general, it would be interesting to explore the central problem of self-modification safety under different agent and environment models and with different assumptions.

\smallskip
\noindent
\textbf{Bounded rationality models.} We analyzed a model of bounded-rationality with a strict upper bound on the size of errors (of several kinds). While this shows that even agents guaranteed to have small errors may self-modify in detrimental ways, the analysis would be significantly different for fully stochastic bounded rationality models (e.g. negligible expected errors with non-negligible variance). One such model of interest is Information-Theoretic Bounded Rationality of \citet{ortega2015informationtheoretic}.

\smallskip
\noindent
\textbf{Awareness of own bounded-rationality.} Whatever underlying decision procedure the agent uses somehow \emph{implicitly} takes its $\epsilon$-optimality into account -- in particular since the assumed $\epsilon$-optimality depends on the behavior of future agent versions. In our formulation, we do not assume the agent to have \emph{explicit} knowledge of its bounded rationality model and $\epsilon$, which would at least intuitively seem useful to know. 

Note, however, that in our framing such explicit knowledge would not be necessarily useful, as any deliberation about it is subject to the same error within $\epsilon$-optimality. Therefore it may be interesting to explore bounded rationality models where the information about own bounded rationality could be explicitly reasoned about (with more precision than e.g. modelling the trajectory of the full environment). Would the agent then be more reluctant to self-modify?

\smallskip
\noindent
\textbf{Time horizons and discounting.} Avoiding temporal discounting would likely yield results with stronger implications (Section~\ref{sec:assumptions}). We propose analysing the finite-time undiscounted case, as well as exploring other means of future value aggregation (finite or infinite, as explored by \citet{Bostrom2011}). 

\smallskip
\noindent
\textbf{Model of self-modification.} As noted above, embedded agents with a strong influence on the environment may self-modify by exploiting the environment. However, the extent of this self-modification, and the strength and stability of mechanisms against self-modification (e.g. via modification-dependent utility function) require further research.

\smallskip
\noindent
\textbf{Probabilistic analysis.} Build stochastic models of agent rationality and self-modification, and perform full probabilistic analysis. This may e.g. inform us about required safety margins. In particular, approaches based on statistical physics and information theory seems to be promising here and have already proven fruitful in analyzing existing optimization problems and algorithms.
